The system architecture is based on the CAGE-2 enterprise topology \cite{kiely2023autonomous} and utilizes RAMPART, a simulator to train and evaluate autonomous cyber defense agents in realistic enterprise and OT-integrated network environments.
It is deployed on an OpenStack cloud environment and consists of 15 virtual machine instances: one high-capacity Defender node and fourteen additional hosts distributed across user, enterprise, and operational subnets. 
The enterprise subnet contains mission-critical servers (enterprise0, enterprise1, enterprise2), the user subnet contains five user workstations (user0 – user4) representing typical entry points for adversarial compromise, and the operational subnet hosts OT-integrated nodes (op\_host0 - op\_host2 and op\_server0) running power grid and chemical plant control workloads. Gateway nodes interconnect these subnets, enforcing routed communication paths and enabling realistic lateral movement across network boundaries. 
A separate control network provides management-plane connectivity for orchestration, monitoring, and persistent Velociraptor communication between all endpoints and the Defender. 
The Defender node, provisioned with elevated compute resources, hosts the Velociraptor server and serves as the execution environment for the blue team actions, telemetry aggregation, logging, and reward evaluation. 
The entire topology is instantiated using OpenStack services including Nova for compute provisioning, Neutron for virtual networking and routing, and HEAT templates for declarative orchestration, enabling reproducible deployment of vulnerable services and controlled network communications. 
This emulated infrastructure preserves realistic TCP/IP routing, service discovery, exploit execution, and defensive response mechanisms, thereby providing a high-fidelity hybrid IT-OT environment for evaluating autonomous and human-in-the-loop cyber defense within RAMPART.